\documentclass[11pt,a4paper]{article}

% ============================================================
% Milestone 1 -- Project Proposal & Literature Review
% DIP Team Project (follows the course Section 6.1 template)
% ============================================================
\usepackage[margin=1in]{geometry}
\usepackage[T1]{fontenc}
\usepackage[utf8]{inputenc}
\usepackage{lmodern}
\usepackage{amsmath,amssymb}
\usepackage{graphicx}
\usepackage{booktabs}
\usepackage{array}
\usepackage[round,authoryear]{natbib}
\usepackage{hyperref}
\hypersetup{colorlinks=true, linkcolor=blue!50!black, citecolor=blue!50!black, urlcolor=blue!50!black}
\usepackage{enumitem}
\usepackage{xcolor}
\usepackage{tabularx}
\usepackage{float}


\title{\LARGE Project Proposal \& Literature Review \\[4pt]
\large Lightweight Real-Time Super-Resolution for Edge/Laptop Deployment}

\author{
  Abdelhakeem Osama Abdelhakeem (2600322) \and
  Student B  \and
  Student C 
}

\date{Milestone 1 --- 2 November 2026 \\{} CSE\,681 --- Advanced Digital Image and Video Processing}

\begin{document}
\maketitle

% ============================================================
\begin{abstract}
Single-image super-resolution (SISR) aims to reconstruct a high-resolution
(HR) image from a low-resolution observation. Although modern
deep-learning-based SISR methods achieve strong reconstruction quality,
many require substantial computational resources, making them unsuitable
for CPU-based or resource-constrained deployment.

This project investigates the trade-off between reconstruction quality and
computational efficiency for lightweight real-time super-resolution,
adapting efficiency-oriented evaluation principles from the NTIRE and AIM
challenges \citep{ren2025ntire,aim2025}. A compact $\times2/\times4$
network is proposed, using lightweight residual blocks based on
depthwise-separable convolutions, PixelShuffle upsampling, and an optional
global bicubic skip connection.

The model will be compared against bicubic interpolation and the
lightweight learned baselines FSRCNN and ESPCN
\citep{dong2016fsrcnn,shi2016espcn}, trained on DIV2K
\citep{agustsson2017div2k} and tested on Set5, Set14, BSD100, Urban100 and
Manga109. We report Y-channel PSNR/SSIM, parameters, GMACs and measured
CPU latency, to determine whether a low computational cost also gives
genuinely low real-world CPU latency, rather than relying on FLOP/GMAC
counts alone.
\end{abstract}

% ============================================================
\section{Problem Statement \& Motivation}
\label{sec:motivation}

Single-image super-resolution is an inverse problem in which an HR image is
reconstructed from a lower-resolution (LR) observation. The task connects
directly to the course material on interpolation and resampling, spatial
filtering and convolution, and learned image restoration. However, achieving
high reconstruction quality can require substantial computational resources,
which makes super-resolution models difficult to deploy on CPUs or
resource-constrained devices. This project therefore focuses on the
relationship between image quality and computational efficiency.

The main research question is:

\begin{quote}
\textit{Can a lightweight super-resolution model achieve competitive image
quality while maintaining genuinely low CPU inference latency?}
\end{quote}

The project compares image quality, parameter count, computational cost,
and measured inference time. In particular, it investigates whether a lower
GMAC count necessarily corresponds to lower real-world CPU latency.

% ============================================================
\section{Related Work}
\label{sec:related}

\begin{table}[h]
\centering
\caption{Key references.}
\label{tab:refs}
\small
\begin{tabular}{@{}p{0.30\linewidth}p{0.08\linewidth}p{0.55\linewidth}@{}}
\toprule
\textbf{Reference} & \textbf{Year} & \textbf{Relevance to this project} \\
\midrule
NTIRE 2025 Efficient SR report \citep{ren2025ntire} & 2025 &
Efficiency-oriented evaluation methodology combining quality with
parameters, FLOPs and runtime; adapted to our CPU-oriented
$\times2/\times4$ evaluation. \\

AIM 2025 efficient perceptual SR \citep{aim2025} & 2025 &
Reports quality together with parameters and computational cost; motivates
our quality--efficiency reporting. \\


FSRCNN \citep{dong2016fsrcnn} & 2016 &
Compact fast SR network; our main learned baseline. \\

ESPCN \citep{shi2016espcn} & 2016 &
Introduces sub-pixel convolution (PixelShuffle), used in our model; second
baseline. \\

DIV2K \citep{agustsson2017div2k} & 2017 &
Training and validation dataset. \\
\bottomrule
\end{tabular}
\end{table}

Recent super-resolution research emphasizes not only reconstruction quality
but also computational efficiency and practical deployment constraints.
The NTIRE 2025 Efficient Super-Resolution Challenge investigated efficient
SR using image quality together with computational cost and runtime
\citep{ren2025ntire}. The AIM 2025 efficient perceptual SR benchmark
reports model parameters and computational cost alongside perceptual
metrics \citep{aim2025}. We adopt this practice of reporting quality and
cost together, while our project focuses on PSNR/SSIM and measured CPU
latency.

We adapt the NTIRE 2025 evaluation methodology rather than reproducing its
exact competition setup. In particular, our study evaluates $\times2$ and
$\times4$ super-resolution, measures inference latency on a CPU, and
reports GMACs at a $1280\times720$ output resolution. This differs from the
NTIRE 2025 setup, which focuses on $\times4$ SR and reports computational
cost and runtime under its own evaluation protocol.

FSRCNN introduced a compact convolutional architecture designed for fast SR,
showing that efficient networks can reach competitive quality with far
fewer computations than earlier approaches \citep{dong2016fsrcnn}. ESPCN
proposed sub-pixel convolution, where upscaling is performed at the end of
the network using PixelShuffle, so that most feature extraction happens in
LR space \citep{shi2016espcn}. FSRCNN introduced a compact convolutional architecture designed for fast SR,
showing that efficient networks can reach competitive quality with far fewer
computations than earlier approaches \citep{dong2016fsrcnn}. ESPCN proposed
sub-pixel convolution, where upscaling is performed at the end of the network
using PixelShuffle, so that most feature extraction happens in LR space
\citep{shi2016espcn}. These ideas motivate our use of LR-domain feature
extraction and PixelShuffle-based upsampling. DIV2K is used as the main
training and validation dataset \citep{agustsson2017div2k}.
DIV2K is a high-quality benchmark widely used for training and evaluating
SISR models \citep{agustsson2017div2k}.

Building on these ideas, this project focuses on a reproducible,
CPU-oriented comparison between lightweight architectures. Rather than
attempting state-of-the-art quality, it studies the practical relationship
between model size, GMACs, image quality and measured CPU latency, together
with width/depth ablations and an optional global bicubic skip connection.

% ============================================================
\section{Dataset \& Baseline Plan}
\label{sec:data}

\subsection{Datasets}

Training uses DIV2K (800 HR training images; 100 validation images). LR
images are generated by bicubic downsampling (PIL bicubic resampling).
Training operates on randomly sampled, scale-aligned $96\times96$ HR
patches, with about 20 patches per image per epoch, and augmentation by
horizontal/vertical flips and $90^\circ$ rotation. During development a
validation subset of 20--50 images is used to reduce evaluation time; final
experiments use the complete validation/test protocol where practical.

\begin{table}[h]
\centering
\caption{Dataset summary.}
\label{tab:data}
\small
\begin{tabular}{@{}p{0.14\linewidth}p{0.27\linewidth}p{0.27\linewidth}p{0.22\linewidth}@{}}
\toprule
\textbf{Dataset} & \textbf{Subset used} & \textbf{License} & \textbf{Link} \\
\midrule
DIV2K & 800 train / 20--50 val (dev), 100 val (final) &
Academic research use &
\url{data.vision.ee.ethz.ch/cvl/DIV2K} \\
Set5 & 5 images (test) & Research use & BasicSR dataset guide \\
Set14 & 14 images (test) & Research use & BasicSR dataset guide \\
BSD100 & 100 images (test) & Research use & BasicSR dataset guide \\
Urban100 & 100 images (test) & Research use & BasicSR dataset guide \\
Manga109 & 109 images (test) &
Academic use only; images not redistributed (not placed in the public
repository) & BasicSR dataset guide \\
\bottomrule
\end{tabular}
\end{table}

The BasicSR dataset-preparation guide
(\url{github.com/XPixelGroup/BasicSR/blob/master/docs/DatasetPreparation.md})
is used for the benchmark sets; the final report will document the exact
source and license of each.

\subsection{Baselines}

\begin{enumerate}
    \item \textbf{Bicubic interpolation} as a non-learning baseline.
    \item \textbf{FSRCNN} \citep{dong2016fsrcnn}: main configuration
    $d=56$, $s=12$, $m=4$, retrained on DIV2K at $\times2$ and $\times4$.
    \item \textbf{ESPCN} \citep{shi2016espcn} as an additional lightweight
    learned baseline.
\end{enumerate}

All learned models use identical preprocessing and evaluation: Y-channel
(ITU-R BT.601) PSNR/SSIM with border shaving and 8-bit rounding. CPU
measurements use the same output resolution, hardware, thread count and
timing protocol for all models.

% ============================================================
\section{Proposed Method}
\label{sec:method}

The proposed model, \textbf{EfficientSR}, performs most feature extraction
in the LR domain and postpones spatial upsampling to the end of the network:

\begin{align*}
\text{LR}
&\rightarrow
3{\times}3\text{ Head}
\rightarrow
N\text{ Lightweight Residual Blocks (+ long skip from head)} \\
&\rightarrow
\text{Tail Conv}
\rightarrow
\text{PixelShuffle}
\rightarrow
+\,\text{Bicubic Skip}
\rightarrow
\text{SR}.
\end{align*}

Each residual block consists of a depthwise $3\times3$ convolution, a
pointwise $1\times1$ convolution, a PReLU activation, a second pointwise
$1\times1$ convolution, and a residual connection. The tail convolution
produces $3s^2$ channels ($s$ = scale factor), rearranged by PixelShuffle
into the HR image. The global bicubic skip provides a strong low-cost base
reconstruction:
\[
I_{\mathrm{SR}} = I_{\mathrm{network}} + I_{\mathrm{bicubic}},
\]
so the convolutional network mainly learns the missing high-frequency
detail. An ablation without this skip measures its contribution.

The controllable parameters are the feature width $C$ and the number of
blocks $N$, allowing models with different computational costs to be
evaluated. EfficientSR is developed under an explicit computational budget
of at most $100{,}000$ trainable parameters and at most $10$ GMACs for a
$1280\times720$ output. The expected configurations contain roughly
$2\times10^4$--$6\times10^4$ parameters and about 3--9 GMACs, depending on
scale, width and depth.

Training uses the $L_1$ loss
\[
\mathcal{L} = \left\| I_{\mathrm{SR}} - I_{\mathrm{HR}} \right\|_1,
\]
with the Adam optimizer and a step learning-rate schedule. Because
depthwise convolutions have low theoretical computational cost but may
exhibit different practical behavior on CPUs, measured CPU runtime is
reported together with parameter count and GMACs rather than relying on
theoretical cost alone.

% ============================================================
\section{Evaluation Metrics \& Preliminary Baseline Results}
\label{sec:eval}

\textbf{Image quality:} PSNR and SSIM, computed on the Y channel following
the conventional SR protocol (borders shaved, outputs rounded to 8 bit).

\textbf{Efficiency:} number of trainable parameters; GMACs for a
$1280\times720$ output; measured CPU latency (median of 20 runs, fixed
thread count, after warm-up); and model size where relevant.

The final analysis compares PSNR/SSIM against CPU runtime and computational
cost. Our implementation reports computational complexity in GMACs,
whereas the NTIRE reports use FLOPs. Since one multiply-accumulate
operation is approximately two floating-point operations, FLOPs
$\approx 2\times$ MACs, so GMAC and GFLOP values must not be compared
directly without this conversion.

Table~\ref{tab:prelim} gives the baseline numbers obtained with our
evaluation pipeline for $\times2$. The bicubic row is final for this
milestone. The FSRCNN row is a \emph{preliminary} result from the best
checkpoint of a training run that was still in progress (epoch 48 of 65).
LR images are generated with PIL bicubic resampling rather than MATLAB
\texttt{imresize}, so absolute values may differ slightly from published
tables; all methods here share the same protocol, so comparisons between
them remain fair. The preliminary FSRCNN scores are also below published
values, which we attribute to the short, still unfinished training schedule
and the RGB-input, DIV2K-only setup; we will check this against the final
checkpoint.

\begin{table}[h]
\centering
\caption{Preliminary $\times2$ baseline results: Y-channel PSNR (dB) / SSIM.}
\label{tab:prelim}
\footnotesize
\setlength{\tabcolsep}{4pt}
\begin{tabular}{@{}lcccccc@{}}
\toprule
\textbf{Model} & \textbf{Set5} & \textbf{Set14} & \textbf{BSD100} &
\textbf{Urban100} & \textbf{Manga109} & \textbf{Avg.} \\
\midrule
Bicubic & 33.97 / 0.9330 & 30.56 / 0.8750 & 29.74 / 0.8495 &
27.07 / 0.8456 & 31.25 / 0.9384 & 30.52 / 0.8883 \\
FSRCNN (prelim.) & 35.08 / 0.9441 & 31.50 / 0.8967 & 30.62 / 0.8786 &
27.99 / 0.8702 & 32.71 / 0.9493 & 31.58 / 0.9078 \\
\bottomrule
\end{tabular}
\end{table}

FSRCNN (24,683 parameters, 5.607 GMACs for a $1280\times720$ output)
improves the average PSNR by about 1.06\,dB over bicubic. Its measured
end-to-end CPU time (4 threads on the Intel Core i5-11400H, median) was about
203\,ms, versus about 25\,ms for bicubic. These timings were taken while
training was running on the same machine, so they will be re-measured on an
idle machine for the final report.

% ============================================================
\section{Feasibility \& Compute Plan}
\label{sec:feasibility}

The project is CPU-OK / GPU-light: training and evaluation use small
patches, compact models, and CPU inference is the deployment target.
EfficientSR is developed under an explicit computational budget of at most
$100{,}000$ trainable parameters and at most $10$ GMACs for a
$1280\times720$ output.

Training was run on an NVIDIA GeForce RTX 3060 Laptop GPU (6\,GB), with an
11th Gen Intel Core i5-11400H CPU (6 cores, 12 threads) and 16\,GB of RAM,
with batch size 32 and 50 validation images, using the FSRCNN $\times2$
configuration. With $96\times96$ HR patches and about 20 patches per image,
one epoch took about 512--600\,s for training and 523--609\,s including
validation (roughly 9--10 minutes). The model reached a validation PSNR of
33.84--33.85\,dB (SSIM 0.9275) on the DIV2K validation subset. At this rate
a full 65-epoch run takes about 10 hours. The speed is limited mainly by
decoding large DIV2K images for patch sampling rather than by the network
itself, so cached pre-extracted crops, fewer patches per image, or smaller
models will be used for ablation experiments. Checkpointing and resume
functionality are implemented. Inference is measured on the CPU with a fixed
thread configuration, reporting the median of about 20 runs.



\subsection{Minimum Viable Deliverable (MVP)}

\begin{enumerate}
    \item A complete $\times2$ super-resolution pipeline.
    \item At least one complete $\times4$ experiment, with a matching
    FSRCNN $\times4$ baseline.
    \item Bicubic baseline numbers (PSNR/SSIM) on the test sets
    (Table~\ref{tab:prelim}).
    \item Trained FSRCNN baseline.
    \item At least one trained EfficientSR configuration.
    \item Evaluation on at least one standard benchmark dataset.
    \item PSNR/SSIM, parameter, GMAC and CPU-time measurements.
    \item A recorded demo ($\le$5 minutes) of the trained model enlarging
    images on the CPU.
\end{enumerate}

\subsection{Stretch Goal}

\begin{itemize}
    \item Full $\times2$ and $\times4$ experiments on all five benchmark sets.
    \item Width/depth computational-budget ablation.
    \item Bicubic global-skip ablation.
    \item ESPCN baseline.
    \item A live real-time or near-real-time demonstration.
\end{itemize}

% ============================================================
\section{Risk Assessment}
\label{sec:risk}

\begin{table}[h]
\centering
\caption{Risk assessment.}
\label{tab:risks}
\small
\begin{tabular}{@{}p{0.38\linewidth}p{0.55\linewidth}@{}}
\toprule
\textbf{Risk} & \textbf{Mitigation} \\
\midrule
Slow training (about 9--10 minutes per epoch in our measurements) &
Cached crops, fewer patches per image, checkpointing/resume, limited
validation during development. \\

Depthwise convolution has low GMACs but high CPU latency &
Measure real CPU runtime; compare with a standard-convolution variant at
matched GMACs. \\

Dataset format or preprocessing inconsistencies &
Consistent preprocessing pipeline; document the exact evaluation
protocol. \\

EfficientSR does not outperform FSRCNN at a comparable budget &
Treat the comparison itself as a result and analyze the quality/latency
trade-off. \\

Training instability &
Checkpointing, fixed validation procedure, conservative training
configuration before larger experiments. \\
\bottomrule
\end{tabular}
\end{table}

% ============================================================
\section{Work Breakdown}
\label{sec:workbreakdown}

\begin{table}[H]
\centering
\caption{Milestone task allocation.}
\label{tab:tasks}
\small
\setlength{\tabcolsep}{3pt}
\renewcommand{\arraystretch}{1.15}

\begin{tabularx}{\textwidth}{@{}p{0.10\textwidth}X X X@{}}
\toprule
\textbf{Milestone} &
\textbf{Student A --- Data \& Baseline} &
\textbf{Student B --- Model/Algorithm} &
\textbf{Student C --- Evaluation} \\
\midrule

M1 (Wk 3) &
DIV2K + 5 benchmark-set pipeline; FSRCNN baseline scoped &
Design EfficientSR architecture within the parameter/GMAC budget &
Define PSNR/SSIM + CPU-runtime protocol; bicubic baseline numbers \\

M2 (Wk 6) &
FSRCNN baseline trained; evaluated on test sets &
EfficientSR first checkpoint &
Preliminary PSNR/SSIM-vs-runtime comparison; verify timing pipeline
(1 benchmark set) \\

M3 (Wk 9) &
Final FSRCNN numbers, all test sets &
Final model + parameter/GMAC-budget and bicubic-skip ablations; at least one
$\times4$ EfficientSR experiment &
Full PSNR/SSIM/runtime table across all 5 sets; final report, repo, demo \\

\bottomrule
\end{tabularx}
\end{table}

\noindent
\textbf{Deliverables:} M1: proposal + bicubic baseline numbers; M2: progress
report with one benchmark-set result; M3: final report, repository and demo.
% ============================================================
\section{Expected Outcome}

A reproducible experimental study of lightweight single-image
super-resolution under practical CPU deployment constraints. The project
will produce quantitative comparisons between reconstruction quality and
computational efficiency, and test whether parameter count and GMACs are
sufficient indicators of practical inference speed, or whether choices such
as depthwise convolution create a mismatch between theoretical cost and
measured CPU latency.

% ============================================================
\section{AI Use Disclosure}

AI assistance (ChatGPT and Claude) was used for LaTeX drafting,
clarification of technical concepts, and code debugging. All technical
decisions, experiments, results, and final content were reviewed and
verified by the team.

% ============================================================
\bibliographystyle{unsrtnat}
\bibliography{references}

\end{document}
